\documentclass[11pt,a4paper]{article}
\usepackage[utf8]{inputenc}
\usepackage[T1]{fontenc}
\usepackage{amsmath,amsfonts,amssymb}
\usepackage{graphicx}
\usepackage{hyperref}
\usepackage{listings}
\usepackage{color}
\usepackage{booktabs}
\usepackage{float}
\usepackage{geometry}
\geometry{margin=1in}
\definecolor{zenblue}{RGB}{41,121,255}
\definecolor{zengreen}{RGB}{52,199,89}
\definecolor{codegray}{RGB}{240,240,240}
\hypersetup{colorlinks=true,linkcolor=zenblue,urlcolor=zenblue,citecolor=zenblue}

\lstset{
  backgroundcolor=\color{codegray},
  basicstyle=\ttfamily\small,
  breaklines=true,
  frame=single
}

\title{\textbf{Zen-Coder-Flash: A Distilled Qwen-Coder Variant for Low-Latency Code Completion}\\
\large Technical Report v2025.02}
\author{Antje Worring, Zach Kelling \\ Zen LM Research Team\\
\texttt{research@zenlm.org}}
\date{February 2025}

\begin{document}
\maketitle

\begin{abstract}
\textbf{Zen-Coder-Flash} is a smaller code-completion student derived by knowledge distillation
from Alibaba's open-weight \textbf{Qwen2.5-Coder} \cite{qwen25coder} (Apache-2.0). The teacher is the
Qwen2.5-Coder-14B checkpoint and the student is initialized from a smaller Qwen2.5-Coder checkpoint
(7B-class); we do \emph{not} pretrain a new model from scratch and we do \emph{not} introduce a novel
``Zen MoDE'' architecture---both the teacher and the student backbone are Qwen's, the tokenizer and
fill-in-the-middle (FIM) format are Qwen's, and our contribution is the distillation recipe and the
deployment packaging. This report describes that recipe: sequence-level KD with intermediate-layer
feature alignment, the latency-oriented serving configuration, and the optional use of the student as
a draft model for speculative decoding against the larger teacher. Benchmark and latency figures
quoted here are illustrative \emph{targets} for the recipe, not validated leaderboard claims; for
authoritative quality numbers on the underlying models we refer the reader to the Qwen2.5-Coder
technical report \cite{qwen25coder}.
\end{abstract}

\tableofcontents
\newpage

%% ─────────────────────────────────────────────────────────────────────────────
\section{Introduction}

IDE code completion operates under strict latency constraints. Empirical studies of developer
experience show that completion latency above 150ms breaks the perceived ``inline'' feeling and
causes developers to wait or dismiss suggestions \cite{svyatkovskiy2021fast}. At P95, this
requires generating a 10--30 token completion in well under 100ms, leaving approximately 30--50ms
budget for network round-trip, model inference, and post-processing.

General-purpose code models at 13--14B parameters are larger than necessary for this budget on widely
available server hardware. Zen-Coder-Flash addresses this by distilling Qwen2.5-Coder-14B
\cite{qwen25coder} into a smaller ($\sim$7--8B-class) Qwen2.5-Coder student, aiming to preserve code
quality while reducing serving latency.

Scope of contributions (engineering recipe, not new science):

\begin{itemize}
  \item A distillation recipe applying sequence-level KD with intermediate-layer feature alignment to
        a Qwen2.5-Coder teacher/student pair \cite{qwen25coder}.
  \item A latency-oriented serving configuration targeting sub-50ms P95 time-to-first-token for short
        autocomplete prefixes on a single A100-class GPU.
  \item Optional speculative decoding: the distilled student serves as a draft model for the larger
        Qwen2.5-Coder-14B teacher, a standard lossless-speedup setup
        \cite{leviathan2023speculative,chen2023accelerating}.
\end{itemize}

\noindent The numeric targets in the tables below (e.g.\ P95 latency, throughput, pass@1) are design
goals for this recipe on the stated hardware, \emph{not} independently audited benchmark results.
Authoritative quality numbers for the Qwen2.5-Coder family are in the upstream report
\cite{qwen25coder}; we do not restate them as our own.

%% ─────────────────────────────────────────────────────────────────────────────
\section{Architecture}

\subsection{Student Model Specification}

Both the teacher and the student are Qwen2.5-Coder \cite{qwen25coder} checkpoints; the student is a
smaller member of the same family (7B-class), not a bespoke ``Zen MoDE'' architecture. The
hyperparameters below are those of the corresponding Qwen2.5-Coder checkpoints and are reproduced from
Qwen's model cards; the depth-to-width ratio and the shared vocabulary/RoPE settings follow from using
the same upstream family for teacher and student.

\begin{table}[H]
\centering
\caption{Student and teacher hyperparameters (both Qwen2.5-Coder \cite{qwen25coder}).}
\label{tab:arch}
\begin{tabular}{lcc}
\toprule
\textbf{Hyperparameter} & \textbf{Student (Qwen2.5-Coder-7B)} & \textbf{Teacher (Qwen2.5-Coder-14B)} \\
\midrule
Parameters (total)           & $\sim$7.6B & $\sim$14.7B \\
Layers                       & 28 & 48 \\
Hidden dimension             & 3{,}584 & 5{,}120 \\
KV heads (GQA)               & 4  & 8 \\
Vocabulary size              & 151{,}936 & 151{,}936 \\
Context length               & 32{,}768 & 32{,}768 \\
Position encoding            & RoPE ($\theta = 1{,}000{,}000$) & RoPE ($\theta = 1{,}000{,}000$) \\
\bottomrule
\end{tabular}
\end{table}

\noindent (Figures are the upstream Qwen2.5-Coder specifications; the exact student checkpoint chosen
for a given deployment may differ in size, but it is in all cases an existing Qwen2.5-Coder model, not
a newly designed architecture.)

\subsection{Architectural Choices for Latency}

The student's smaller hidden dimension and lower layer count are the latency drivers: matrix-multiply
cost per layer scales with hidden dimension, and total cost scales with layer count. Choosing a
smaller existing Qwen2.5-Coder checkpoint as the student is what buys the latency reduction; we do not
re-architect the model.

Grouped-query attention (inherited from the Qwen2.5-Coder backbone) keeps KV-cache requirements modest
and enables large effective batch sizes. As an order-of-magnitude illustration of KV-cache size for an
8-KV-head, 128-dim-per-head configuration:

\begin{equation}
  \text{KV cache (per layer)} = 2 \cdot n_{\text{kv}} \cdot d_k \cdot S \cdot 2 = 2 \cdot 8 \cdot 128 \cdot 32{,}768 \cdot 2 \approx 134 \text{ MB}
\end{equation}

A modest per-layer KV-cache footprint at 32K context leaves headroom for the BF16 weights of a
7--8B-class student on a single 24 GB consumer GPU.

%% ─────────────────────────────────────────────────────────────────────────────
\section{Distillation Methodology}

\subsection{Overview}

Because both teacher and student are existing Qwen2.5-Coder \cite{qwen25coder} checkpoints, there is
\textbf{no from-scratch pretraining phase}. The recipe is two steps applied on top of the upstream
weights:

\begin{enumerate}
  \item \textbf{Knowledge distillation}: starting from the pretrained Qwen2.5-Coder student checkpoint,
        distill from the Qwen2.5-Coder-14B teacher using combined token-level KD and
        intermediate-feature-alignment losses.
  \item \textbf{Instruction fine-tuning}: a light SFT pass on instruction-code pairs to restore
        instruction-following after the KD pass.
\end{enumerate}

\noindent Earlier drafts described a ``Zen MoDE distillation pipeline'' and a 500B-token from-scratch
warm-up over a proprietary 2.5T corpus. Neither is accurate: the student is not trained from scratch
(it is an existing Qwen2.5-Coder checkpoint) and there is no proprietary multi-trillion-token corpus
behind this work. Those claims have been removed.

\subsection{Initialization}

The student is \emph{initialized from a pretrained Qwen2.5-Coder checkpoint}, not randomly initialized.
This is precisely why no warm-up pretraining is needed: the student already has strong code priors from
Qwen's upstream pretraining, so distillation begins from a competent starting point rather than from
poor initial approximations of the teacher.

\paragraph{Distillation-run settings (illustrative).}
\begin{itemize}
  \item Optimizer: AdamW, LR $1 \times 10^{-4}$, cosine to $5 \times 10^{-6}$
  \item Batch: 2M tokens
  \item FIM rate: 50\% (Qwen2.5-Coder FIM format)
  \item Scale: tens of billions of distillation tokens (orders of magnitude below a pretraining run)
\end{itemize}

\subsection{Knowledge Distillation}

The distillation loss combines three terms:

\begin{equation}
  \mathcal{L}_{\text{KD}} = \alpha \mathcal{L}_{\text{CE}} + \beta \mathcal{L}_{\text{KL}} + \gamma \mathcal{L}_{\text{feat}}
\end{equation}

\paragraph{Cross-entropy loss ($\mathcal{L}_{\text{CE}}$).}
Standard next-token prediction on ground-truth labels:

\begin{equation}
  \mathcal{L}_{\text{CE}} = -\sum_t \log p_s(y_t \mid y_{<t}, x)
\end{equation}

\paragraph{KL divergence loss ($\mathcal{L}_{\text{KL}}$).}
Token-level KL divergence between student and teacher output distributions at temperature $T$:

\begin{equation}
  \mathcal{L}_{\text{KL}} = \sum_t D_{\text{KL}}\!\left(p_t\!\left(\cdot; T\right) \| p_s\!\left(\cdot; T\right)\right)
\end{equation}

with temperature $T = 1.0$ (we find $T > 1$ unnecessarily smooths code-specific distributions
where the correct token is often highly peaked). Top-$k$ filtering at $k=50$ is applied to
the teacher logits before KL computation to reduce noise from near-zero probabilities.

\paragraph{Intermediate feature alignment ($\mathcal{L}_{\text{feat}}$).}
We align intermediate hidden states between the student and teacher at evenly spaced layers
using a learned linear projection $W_{\text{proj}}$ mapping the student's hidden dimension
($d_s = 3584$ for the 7B student) to the teacher's ($d_t = 5120$):

\begin{equation}
  \mathcal{L}_{\text{feat}} = \frac{1}{|M|} \sum_{(i,j) \in M} \left\| h^s_i W_{\text{proj}} - h^t_j \right\|_2^2
\end{equation}

where $M$ maps a set of evenly spaced student layers to evenly spaced teacher layers (the student has
fewer layers than the teacher, so the mapping is many-to-fewer). The projection matrix $W_{\text{proj}}$
is trained jointly with the student.

Loss coefficients: $\alpha = 0.3$, $\beta = 0.5$, $\gamma = 0.2$.

\begin{table}[H]
\centering
\caption{Distillation phase hyperparameters.}
\label{tab:kd}
\begin{tabular}{lc}
\toprule
\textbf{Parameter} & \textbf{Value} \\
\midrule
Learning rate      & $1 \times 10^{-4}$ $\to$ $5 \times 10^{-6}$ (cosine) \\
Batch size         & 2M tokens \\
KD temperature     & 1.0 \\
Top-$k$ filtering  & 50 \\
CE coefficient $\alpha$ & 0.3 \\
KL coefficient $\beta$  & 0.5 \\
Feature coefficient $\gamma$ & 0.2 \\
Scale              & distillation-scale (tens of billions of tokens) \\
\bottomrule
\end{tabular}
\end{table}

\subsection{Speculative Decoding Setup}

The distilled student can serve as a draft model in a speculative decoding pipeline with the
Qwen2.5-Coder-14B teacher \cite{qwen25coder} as the verifier. The draft model generates $\gamma = 6$
speculative tokens per step; the verifier accepts or rejects in parallel.

Expected speedup with acceptance rate $\alpha$:

\begin{equation}
  \text{Speedup} = \frac{1 + \alpha + \alpha^2 + \cdots + \alpha^\gamma}{\text{(1 verifier step cost)} / \text{(draft step cost)}}
\end{equation}

A typical draft-model acceptance rate for code completion is in the 0.7--0.85 range; at the upper end
this corresponds to a single-host wall-clock speedup of roughly $2$--$3\times$ over running the teacher
alone, consistent with the speculative-decoding literature \cite{leviathan2023speculative,chen2023accelerating}.
The exact figure is workload-dependent and should be measured per deployment.

%% ─────────────────────────────────────────────────────────────────────────────
\section{Expected Performance (Targets)}

\textbf{The numbers in this section are design targets for the recipe, not audited benchmark results.}
We have not run an independent leaderboard evaluation of the distilled student, and we do not claim to
beat the upstream Qwen2.5-Coder models. For authoritative code-quality numbers on the teacher and on
same-size Qwen2.5-Coder checkpoints, see the Qwen2.5-Coder technical report \cite{qwen25coder}; those
upstream numbers bound what a distilled student can be expected to recover.

\subsection{Code Quality (Target Recovery)}

The aim of distillation is to recover most of the teacher's code-quality on the smaller student. As a
rough target we aim for $\sim$90\%+ recovery of the teacher's pass@1 on HumanEval/MBPP; the student
cannot exceed the teacher, and a same-size Qwen2.5-Coder checkpoint is the relevant floor. We report no
fabricated head-to-head table here; per-deployment numbers should be measured against the upstream
Qwen2.5-Coder baselines \cite{qwen25coder}.

\subsection{Latency (Target)}

The serving goal is sub-50ms P95 time-to-first-token for short ($\sim$512-token) autocomplete prefixes
on a single A100-class GPU in BF16 with a paged-attention server \cite{kwon2023efficient}, which a
7--8B-class model can plausibly meet where the 14B teacher cannot. These are engineering targets to be
validated on the deployment hardware, not measured results reproduced here.

\subsection{Speculative Decoding (Expected)}

Using the distilled student as a draft model for the Qwen2.5-Coder-14B teacher is expected to yield a
$2$--$3\times$ wall-clock speedup over the teacher alone at the teacher's exact output distribution,
since speculative decoding is lossless modulo floating-point differences in acceptance sampling
\cite{leviathan2023speculative,chen2023accelerating}. Realized speedup depends on the acceptance rate
for the target workload and must be measured.

%% ─────────────────────────────────────────────────────────────────────────────
\section{Distillation Ablation}

We do not report a table of measured pass@1 values for each loss configuration, as we have not run an
audited evaluation to support specific numbers. Qualitatively, and consistent with the distillation
literature \cite{hinton2015distilling,kim2016seq,gu2024minillm,jiao2020tinybert}, we expect both the
KL-divergence term (token-level distribution matching) and the intermediate-feature-alignment term
(hidden-state matching) to contribute independently over a cross-entropy-only baseline, with the
combined loss outperforming either term alone. Latency is unaffected by the choice of distillation loss
(it changes training, not the served architecture). Any quantitative ablation should be reported with
the deployment's own measured numbers.

%% ─────────────────────────────────────────────────────────────────────────────
\section{IDE Integration and Deployment}

\subsection{Deployment Architecture}

Zen-Coder-Flash is designed for always-on inference deployment:

\begin{lstlisting}[language=bash, caption={Starting a Zen-Coder-Flash inference server.}]
# Single GPU deployment (24GB VRAM minimum)
docker run --gpus=1 \
  -e MODEL=zen-coder-flash-8b \
  -e MAX_BATCH_SIZE=32 \
  -e MAX_CONTEXT=32768 \
  -p 8080:8080 \
  ghcr.io/zenlm/inference:latest

# Speculative decoding with Qwen2.5-Coder-14B verifier (2x A100)
docker run --gpus=2 \
  -e MODEL=Qwen/Qwen2.5-Coder-14B-Instruct \
  -e DRAFT_MODEL=zen-coder-flash \
  -e SPECULATIVE_GAMMA=6 \
  -p 8080:8080 \
  ghcr.io/zenlm/inference:latest
\end{lstlisting}

\subsection{Continue Extension Configuration}

The Continue VS Code extension can be configured to use Zen-Coder-Flash for autocomplete
with the following configuration:

\begin{lstlisting}[language={},caption={Continue config.json for Zen-Coder-Flash.}]
{
  "tabAutocompleteModel": {
    "title": "Zen-Coder-Flash",
    "provider": "openai",
    "model": "zen-coder-flash-8b",
    "apiBase": "https://api.zenlm.org/v1",
    "apiKey": "<your-api-key>"
  },
  "tabAutocompleteOptions": {
    "useCopyBuffer": false,
    "maxPromptTokens": 1024,
    "prefixPercentage": 0.85
  }
}
\end{lstlisting}

%% ─────────────────────────────────────────────────────────────────────────────
\section{Related Work}

Knowledge distillation for language models was introduced by \cite{hinton2015distilling}
and applied to transformers in DistilBERT \cite{sanh2019distilbert}. For generative models,
sequence-level KD \cite{kim2016seq} and token-level KL distillation \cite{gu2024minillm}
have emerged as standard approaches. Intermediate layer alignment follows the TinyBERT
\cite{jiao2020tinybert} and PKD \cite{sun2019patient} patterns. Speculative decoding
\cite{leviathan2023speculative, chen2023accelerating} enables lossless speedup using a small
draft model; our use of the distilled student as a draft for the Qwen2.5-Coder-14B teacher
\cite{qwen25coder} follows this paradigm.

The base models themselves are Qwen2.5-Coder \cite{qwen25coder}; Zen-Coder-Flash contributes a
distillation recipe and a latency-oriented serving configuration on top of those open-weight models,
not a new model. Fast code completion has also been explored in systems such as FauxPilot.

%% ─────────────────────────────────────────────────────────────────────────────
\section{Limitations}

A distilled student is bounded by its teacher: it cannot exceed Qwen2.5-Coder-14B
\cite{qwen25coder}, and the gap typically widens on complex multi-file reasoning tasks (e.g.\
SWE-bench), which benefit disproportionately from the teacher's additional capacity. The student's
context window is bounded by the chosen Qwen2.5-Coder checkpoint. Speculative-decoding speedups are
reduced under the variable-length acceptance patterns typical of code generation. Finally, the quality
and latency figures in this report are \emph{targets}, not audited measurements; deployments should
validate against the upstream Qwen2.5-Coder baselines and their own hardware.

%% ─────────────────────────────────────────────────────────────────────────────
\section{Conclusion}

Zen-Coder-Flash is a smaller, lower-latency code-completion student obtained by knowledge distillation
from Alibaba's open-weight Qwen2.5-Coder-14B \cite{qwen25coder} into a smaller Qwen2.5-Coder checkpoint.
It is not a from-scratch model and does not use a ``Zen MoDE'' architecture; the teacher, the student
backbone, the tokenizer, and the FIM format are all Qwen's, and the work is released under the upstream
Apache-2.0 license with attribution. Our contribution is the distillation recipe (sequence-level KD plus
intermediate-feature alignment), a latency-oriented serving configuration, and the option to use the
student as a speculative-decoding draft for the teacher. The latency and quality figures here are design
targets for that recipe, to be validated per deployment against the upstream Qwen2.5-Coder baselines.

%% ─────────────────────────────────────────────────────────────────────────────
\begin{thebibliography}{99}

\bibitem{qwen25coder}
B.~Hui, J.~Yang, et al.\ (Qwen Team, Alibaba), ``Qwen2.5-Coder Technical Report,''
\textit{arXiv:2409.12186}, 2024. Apache-2.0 (0.5B/1.5B/7B/14B/32B).
\url{https://github.com/QwenLM/Qwen2.5-Coder}.

\bibitem{svyatkovskiy2021fast}
A.~Svyatkovskiy et al., ``Fast and Low-Cost Code Intelligence with Lightweight Models,''
\textit{arXiv:2104.14765}, 2021.

\bibitem{hinton2015distilling}
G.~Hinton, O.~Vinyals, and J.~Dean, ``Distilling the Knowledge in a Neural Network,''
\textit{arXiv:1503.02531}, 2015.

\bibitem{sanh2019distilbert}
V.~Sanh et al., ``DistilBERT, a distilled version of BERT,'' \textit{arXiv:1910.01108}, 2019.

\bibitem{kim2016seq}
Y.~Kim and A.~Rush, ``Sequence-Level Knowledge Distillation,'' \textit{EMNLP}, 2016.

\bibitem{gu2024minillm}
Y.~Gu et al., ``MiniLLM: Knowledge Distillation of Large Language Models,''
\textit{ICLR}, 2024.

\bibitem{jiao2020tinybert}
X.~Jiao et al., ``TinyBERT: Distilling BERT for Natural Language Understanding,''
\textit{EMNLP}, 2020.

\bibitem{sun2019patient}
S.~Sun et al., ``Patient Knowledge Distillation for BERT Compression,''
\textit{EMNLP}, 2019.

\bibitem{leviathan2023speculative}
Y.~Leviathan, M.~Kalman, and Y.~Matias, ``Fast Inference from Transformers via Speculative
Decoding,'' \textit{ICML}, 2023.

\bibitem{chen2023accelerating}
C.~Chen et al., ``Accelerating Large Language Model Decoding with Speculative Sampling,''
\textit{arXiv:2302.01318}, 2023.

\bibitem{kwon2023efficient}
W.~Kwon et al., ``Efficient Memory Management for Large Language Model Serving with PagedAttention,''
\textit{SOSP}, 2023.

\bibitem{bavarian2022fim}
M.~Bavarian et al., ``Efficient Training of Language Models to Fill in the Middle,''
\textit{arXiv:2207.14255}, 2022.

\bibitem{chen2021codex}
M.~Chen et al., ``Evaluating Large Language Models Trained on Code,''
\textit{arXiv:2107.03374}, 2021.

\end{thebibliography}

\end{document}
